\section{Results and Discussion}
\label{sec:results}

All results are reported as five-fold mean best test accuracy \(\pm\) fold-level standard deviation unless otherwise stated.

\subsection{Overall benchmark performance}

The main benchmark asks whether residual topology has a stable winner across datasets and message-passing operators. It does not. Table~\ref{tab:operator_wins} summarizes the full benchmark: across the 24 dataset--operator combinations, MatrixRes is the most frequent winner, with 9 wins and the best mean rank (2.12). MatrixResGated and Plain each win 5 combinations, HorizontalRes wins 3, and VerticalRes wins 2. Thus, matrix-style reuse has the most favorable aggregate profile in this controlled benchmark, but every topology has a setting in which it is the strongest choice.

\begin{table}[!htbp]
\centering
\small
\caption{Winner counts and mean ranks across the 24 dataset--operator combinations. Lower mean rank is better.}
\label{tab:operator_wins}
\begin{tabular}{l c c l}
\toprule
Model & Wins & Mean rank & Interpretation \\
\midrule
Plain & 5 & 3.04 & Strong on selected small or easier settings \\
VerticalRes & 2 & 3.38 & Depth-only reuse helps selectively \\
HorizontalRes & 3 & 3.54 & Competitive when lateral exchange is sufficient \\
MatrixRes & \(\mathbf{9}\) & \(\mathbf{2.12}\) & Most frequent winner in this benchmark \\
MatrixResGated & 5 & 2.92 & Useful when sparse control helps residual traffic \\
\bottomrule
\end{tabular}
\end{table}

\begin{figure}[!htbp]
\centering
\includegraphics[width=0.82\linewidth]{figures/exp/fig_model_win_counts.pdf}
\caption{Model-level winner counts across the 24 dataset--operator combinations.}
\label{fig:model_win_counts}
\end{figure}

Table~\ref{tab:winners_by_operator} gives the cell-level evidence behind the win counts in Table~\ref{tab:operator_wins}. The matrix family is most prominent on MUTAG, AIDS, and Mutagenicity, where MatrixRes or MatrixResGated wins most operator settings. The protein-oriented datasets are more mixed: PROTEINS alternates between Plain and HorizontalRes, while DD alternates among HorizontalRes, MatrixRes, and VerticalRes. ENZYMES remains difficult and does not produce a stable residual-topology preference.

\begin{table}[!htbp]
\centering
\scriptsize
\caption{Winning model for each dataset and message-passing operator. Values in parentheses are mean best test accuracy.}
\label{tab:winners_by_operator}
\resizebox{\linewidth}{!}{%
\begin{tabular}{l c c c c}
\toprule
Dataset & GCNConv & GATConv & SAGEConv & GINConv \\
\midrule
PROTEINS & Plain (0.7044) & HorizontalRes (0.7116) & Plain (0.6999) & HorizontalRes (0.7080) \\
DD & HorizontalRes (0.7181) & MatrixRes (0.7164) & MatrixRes (0.7198) & VerticalRes (0.7224) \\
ENZYMES & MatrixResGated (0.2933) & Plain (0.2717) & VerticalRes (0.2883) & Plain (0.3100) \\
MUTAG & MatrixRes (0.7609) & MatrixResGated (0.7555) & MatrixRes (0.7504) & MatrixResGated (0.8081) \\
AIDS & MatrixRes (0.8365) & MatrixRes (0.8875) & Plain (0.8850) & MatrixResGated (0.9280) \\
Mutagenicity & MatrixRes (0.7784) & MatrixRes (0.7909) & MatrixResGated (0.8031) & MatrixRes (0.8024) \\
\bottomrule
\end{tabular}
}
\end{table}

The ENZYMES outcomes should be interpreted cautiously because the absolute accuracies are low for all compared models. Keeping ENZYMES in the full 24-combination summary preserves the preregistered benchmark scope, but it should not be read as strong evidence for a residual topology by itself. Excluding ENZYMES leaves MatrixRes with 9 wins across 20 dataset--operator combinations and a mean rank of 2.05, so the aggregate MatrixRes preference remains. The split is still context-dependent: MatrixRes is most prominent on the molecular datasets, whereas the protein-oriented datasets produce mixed winners.

\subsection{GCNConv benchmark slice}

To further remove variation caused by message-passing operators and observe the effect of residual topology under a fixed backbone, we report the GCNConv slice separately. Table~\ref{tab:main_benchmark} gives the full slice at the default branch count \(B=3\). Plain is strongest on PROTEINS, HorizontalRes is strongest on DD, MatrixResGated is strongest on ENZYMES, and MatrixRes is strongest on MUTAG, AIDS, and Mutagenicity.

\begin{table}[!htbp]
\centering
\scriptsize
\caption{GCNConv benchmark results with default branch count \(B=3\).}
\label{tab:main_benchmark}
\begin{tabular}{l c c c c c}
\toprule
Dataset & Plain & VerticalRes & HorizontalRes & MatrixRes & MatrixResGated \\
\midrule
PROTEINS & \(\mathbf{0.7044\pm0.0387}\) & \(0.6993\pm0.0340\) & \(0.6992\pm0.0399\) & \(0.6973\pm0.0324\) & \(0.6886\pm0.0397\) \\
DD & \(0.7053\pm0.0517\) & \(0.7143\pm0.0260\) & \(\mathbf{0.7181\pm0.0296}\) & \(0.7127\pm0.0365\) & \(0.7141\pm0.0359\) \\
ENZYMES & \(0.2683\pm0.0343\) & \(0.2550\pm0.0455\) & \(0.2633\pm0.0584\) & \(0.2750\pm0.0577\) & \(\mathbf{0.2933\pm0.0470}\) \\
MUTAG & \(0.7556\pm0.0439\) & \(0.7451\pm0.0695\) & \(0.7343\pm0.0433\) & \(\mathbf{0.7609\pm0.0715}\) & \(0.7556\pm0.0721\) \\
AIDS & \(0.8215\pm0.0046\) & \(0.8210\pm0.0233\) & \(0.8130\pm0.0297\) & \(\mathbf{0.8365\pm0.0412}\) & \(0.8150\pm0.0231\) \\
Mutagenicity & \(0.7644\pm0.0188\) & \(0.7667\pm0.0212\) & \(0.7738\pm0.0134\) & \(\mathbf{0.7784\pm0.0185}\) & \(0.7701\pm0.0165\) \\
\bottomrule
\end{tabular}
\end{table}

\begin{figure}[!htbp]
\centering
\includegraphics[width=0.92\linewidth]{figures/exp/fig_main_benchmark_gcnconv.pdf}
\caption{GCNConv benchmark at \(B=3\). Bars show five-fold mean best test accuracy with fold-level standard-deviation error bars.}
\label{fig:main_benchmark}
\end{figure}

The GCNConv slice supports the same conclusion as the full benchmark: residual topology matters, but the best topology depends on the dataset. MatrixRes is the strongest default choice on three of six GCNConv datasets, while simpler residual paths are still preferable on PROTEINS and DD. The ENZYMES winner is reported for completeness, but its low absolute accuracy limits the strength of that comparison.

\subsection{Synthetic multi-class scaling}

The synthetic benchmark asks a different question from the real-data benchmark: does residual topology behave differently when the structural class granularity is increased under controlled graph-generation rules? Table~\ref{tab:synthetic_class_scaling} summarizes the class-count trend after averaging over the five synthetic graph families and four message-passing operators. The two- and three-class settings do not favor matrix-style reuse: Plain is best at \(C=2\), and HorizontalRes is best at \(C=3\). From \(C=4\) onward, matrix-style reuse becomes more competitive. MatrixResGated is the strongest mean-accuracy model at \(C=4\), \(C=5\), \(C=6\), and \(C=8\), while VerticalRes is strongest at \(C=7\).

\begin{table}[!htbp]
\centering
\scriptsize
\caption{Synthetic multi-class scaling averaged over five graph families and four operators. NAcc denotes chance-normalized accuracy.}
\label{tab:synthetic_class_scaling}
\begin{tabular}{c l c c c c c}
\toprule
Classes & Best model & Best Acc. & Plain Acc. & MatrixRes Acc. & MatrixResGated Acc. & Best NAcc \\
\midrule
2 & Plain & \(\mathbf{0.9585}\) & \(\mathbf{0.9585}\) & 0.9481 & 0.9578 & 0.9170 \\
3 & HorizontalRes & \(\mathbf{0.8800}\) & 0.8774 & 0.8600 & 0.8772 & 0.8200 \\
4 & MatrixResGated & \(\mathbf{0.8111}\) & 0.8002 & 0.8083 & \(\mathbf{0.8111}\) & 0.7481 \\
5 & MatrixResGated & \(\mathbf{0.7296}\) & 0.7233 & 0.7280 & \(\mathbf{0.7296}\) & 0.6621 \\
6 & MatrixResGated & \(\mathbf{0.6878}\) & 0.6719 & 0.6827 & \(\mathbf{0.6878}\) & 0.6254 \\
7 & VerticalRes & \(\mathbf{0.6524}\) & 0.6314 & 0.6429 & 0.6395 & 0.5945 \\
8 & MatrixResGated & \(\mathbf{0.6035}\) & 0.5865 & 0.6002 & \(\mathbf{0.6035}\) & 0.5469 \\
\bottomrule
\end{tabular}
\end{table}

\begin{figure}[!htbp]
\centering
\includegraphics[width=0.92\linewidth]{figures/exp/fig_synthetic_multiclass_scaling.pdf}
\caption{Synthetic multi-class scaling from two to eight classes. Each point averages five synthetic graph families and four message-passing operators.}
\label{fig:synthetic_multiclass_scaling}
\end{figure}

Table~\ref{tab:synthetic_high_class} focuses on the higher-class regime \(C=4,\ldots,8\), where the task requires finer structural discrimination. MatrixResGated has the best mean accuracy and normalized accuracy in this regime, improving accuracy from 0.6827 for Plain to 0.6943. MatrixRes has the strongest macro-F1 among the compared models and is close to MatrixResGated in normalized accuracy. The class-scaling curve in Figure~\ref{fig:synthetic_multiclass_scaling} also shows a stability pattern: from \(C=2\) to \(C=8\), MatrixRes loses 0.3479 accuracy points, compared with 0.3720 for Plain. Thus, the synthetic benchmark supports a narrow but useful interpretation: matrix-style reuse is not necessary for easy low-class settings, but it becomes more favorable as the structural labels become more fine-grained.

\begin{table}[!htbp]
\centering
\small
\caption{Synthetic high-class regime averaged over \(C=4,\ldots,8\), five graph families, and four operators.}
\label{tab:synthetic_high_class}
\begin{tabular}{l c c c}
\toprule
Model & Accuracy & Macro-F1 & NAcc \\
\midrule
Plain & 0.6827 & 0.6703 & 0.6183 \\
VerticalRes & 0.6931 & 0.6834 & 0.6307 \\
HorizontalRes & 0.6846 & 0.6722 & 0.6210 \\
MatrixRes & 0.6924 & \(\mathbf{0.6880}\) & 0.6300 \\
MatrixResGated & \(\mathbf{0.6943}\) & 0.6876 & \(\mathbf{0.6324}\) \\
\bottomrule
\end{tabular}
\end{table}

\subsection{Branch-count ablation}

Branch count is not merely a capacity parameter: it also changes residual-path density, so it requires a separate ablation. Table~\ref{tab:branch_best} reports the best and worst branch counts for each topology. On PROTEINS, the useful region is moderate: HorizontalRes peaks at \(B=4\), MatrixRes peaks at \(B=6\), and MatrixResGated peaks at \(B=3\). On DD, the useful region is smaller: HorizontalRes and MatrixRes peak at \(B=2\), while MatrixResGated is strongest already at \(B=1\). Thus, PROTEINS benefits from moderate branch budgets, whereas DD favors smaller budgets. Branch expansion must be matched to task structure.

\begin{table}[!htbp]
\centering
\small
\caption{Best and worst branch-count settings in the ablation.}
\label{tab:branch_best}
\begin{tabular}{l l c c c c}
\toprule
Dataset & Model & Best \(B\) & Best accuracy & Worst \(B\) & Worst accuracy \\
\midrule
PROTEINS & HorizontalRes & \(\mathbf{4}\) & \(\mathbf{0.7125\pm0.0286}\) & 1 & \(0.6801\pm0.0518\) \\
PROTEINS & MatrixRes & \(\mathbf{6}\) & \(\mathbf{0.7062\pm0.0210}\) & 1 & \(0.6909\pm0.0586\) \\
PROTEINS & MatrixResGated & \(\mathbf{3}\) & \(\mathbf{0.7098\pm0.0219}\) & 4 & \(0.6711\pm0.0471\) \\
DD & HorizontalRes & \(\mathbf{2}\) & \(\mathbf{0.7275\pm0.0304}\) & 4 & \(0.7071\pm0.0302\) \\
DD & MatrixRes & \(\mathbf{2}\) & \(\mathbf{0.7206\pm0.0411}\) & 7 & \(0.7071\pm0.0390\) \\
DD & MatrixResGated & \(\mathbf{1}\) & \(\mathbf{0.7283\pm0.0244}\) & 7 & \(0.7053\pm0.0456\) \\
\bottomrule
\end{tabular}
\end{table}

\begin{figure}[!htbp]
\centering
\includegraphics[width=0.92\linewidth]{figures/exp/fig_branch_count_ablation.pdf}
\caption{Branch-count ablation for HorizontalRes, MatrixRes, and MatrixResGated on PROTEINS and DD under GCNConv. The curves are not monotonic: accuracy improves only within an effective operating region before saturating or declining.}
\label{fig:branch_ablation}
\end{figure}

Table~\ref{tab:branch_best} and Figure~\ref{fig:branch_ablation} show that branch count should not be read as a simple capacity knob. Increasing \(B\) changes the residual graph itself: it adds branches, residual paths, residual traffic, and optimization interactions. PROTEINS tolerates moderate branch expansion, while DD is more sensitive to redundant branch expansion.

\subsection{Mechanism analysis: diversity, similarity, gradients, and residual traffic}

To explain why branch-count changes affect performance rather than merely report accuracy, we analyze branch diversity, branch cosine similarity, branch CKA similarity, representative gradient norms, residual-count totals, active ratios, and gate values. Table~\ref{tab:mechanism_snapshot} reports these mechanism metrics at the best branch count for each ablation target. Diversity measures how far apart branch representations are; branch cosine and branch CKA measure representational redundancy or alignment; gradient norm captures optimization signal; residual count and active ratio describe the amount of residual traffic injected into the model. These measurements are diagnostic correlations rather than causal interventions.

\begin{table}[!htbp]
\centering
\scriptsize
\caption{Mechanism metrics at the best branch count for each ablation target.}
\label{tab:mechanism_snapshot}
\resizebox{\linewidth}{!}{%
\begin{tabular}{l l c c c c c c c c c}
\toprule
Dataset & Model & \(B\) & Acc. & Diversity & Cosine & CKA & Grad. & Residuals & Active & Gate \\
\midrule
PROTEINS & HorizontalRes & 4 & \(\mathbf{0.7125}\) & 1.1450 & 0.8417 & 0.9797 & 0.0157 & 24.0 & 0.2519 & 0.7911 \\
PROTEINS & MatrixRes & 6 & 0.7062 & 3.7755 & 0.9173 & 0.9634 & 0.0182 & 88.0 & 0.4273 & 0.7757 \\
PROTEINS & MatrixResGated & 3 & 0.7098 & 1.3555 & 0.9790 & 0.9988 & 0.0316 & 37.0 & 0.3121 & 0.7833 \\
DD & HorizontalRes & 2 & 0.7275 & 0.1687 & 0.9969 & 1.0000 & 0.0732 & 6.0 & 0.3492 & 0.8375 \\
DD & MatrixRes & 2 & 0.7206 & 0.1834 & 0.9991 & 1.0000 & 0.0796 & 14.0 & 0.5384 & 0.8222 \\
DD & MatrixResGated & 1 & \(\mathbf{0.7283}\) & 0.0000 & 1.0000 & 1.0000 & 0.0821 & 2.0 & 0.2352 & 0.8221 \\
\bottomrule
\end{tabular}
}
\end{table}

The PROTEINS ablation shows that moderate branch separation can be useful. HorizontalRes improves from \(B=1\) to \(B=4\) while branch diversity increases from 0.0000 to 1.1450 and branch cosine decreases from 1.0000 to 0.8417. However, diversity alone does not guarantee better accuracy: at \(B=8\), diversity rises further to 1.5441, but accuracy falls to 0.6881. MatrixRes follows a similar pattern at a larger scale. It peaks at \(B=6\), where branch diversity reaches 3.7755, but adding still more branches does not improve accuracy.

DD is more conservative. HorizontalRes reaches its best result at \(B=2\), where branch diversity is only 0.1687 and branch cosine remains 0.9969. MatrixRes also peaks at \(B=2\), with similarly high branch cosine (0.9991). Larger branch counts increase diversity on DD, but the performance does not improve. This pattern suggests that larger residual neighborhoods are not useful on DD unless the added branches produce sufficiently useful functional separation.

The sparse/gated model clarifies the role of residual control. MatrixResGated has an active ratio of 0.3121 at the best PROTEINS setting and 0.2352 at the best DD setting. These values are consistent with sparse filtering suppressing residual entries rather than behaving as a dense residual path. At the same time, MatrixResGated does not dominate MatrixRes overall, so sparse control should be treated as a useful traffic-control mechanism rather than as a universal improvement.

\begin{figure}[!htbp]
\centering
\includegraphics[width=0.92\linewidth]{figures/exp/fig_mechanism_branch_dynamics.pdf}
\caption{Mechanism summary linking branch-count accuracy to branch diversity, branch cosine similarity, branch CKA, and mean gradient norm.}
\label{fig:mechanism}
\end{figure}

\subsection{MatrixRes and MatrixResGated sensitivity}

The fold-0 sensitivity scan identifies candidate operating regions for MatrixRes and MatrixResGated. These scans are used for diagnosis rather than final claims. On PROTEINS, MatrixRes does not improve over its fold-0 baseline in the tested sweep, while MatrixResGated has a fold-0 sparsity candidate: \(\lambda=0.02\) improves fold-0 accuracy from 0.6726 to 0.6951. Stronger sparsification is harmful. On DD, both MatrixRes and MatrixResGated prefer lighter hidden dimensions in the fold-0 scan: MatrixRes reaches 0.7595 with \(d=32\), and MatrixResGated reaches 0.7511 with \(d=32\). MatrixResGated also reaches 0.7511 with \(lr=0.001\) and \(gate\_init=0.2\).

\begin{table}[!htbp]
\centering
\small
\caption{Best fold-0 sensitivity settings for MatrixRes and MatrixResGated.}
\label{tab:sensitivity}
\begin{tabular}{l l c c l}
\toprule
Dataset & Model & Baseline & Best setting & Best accuracy \\
\midrule
DD & MatrixRes & 0.7468 & \(dim=32\) & 0.7595 \\
DD & MatrixResGated & 0.7300 & \(dim=32\) & 0.7511 \\
PROTEINS & MatrixRes & 0.6637 & baseline & 0.6637 \\
PROTEINS & MatrixResGated & 0.6726 & \(\lambda=0.02\) & 0.6951 \\
\bottomrule
\end{tabular}
\end{table}

\begin{figure}[!htbp]
\centering
\includegraphics[width=0.92\linewidth]{figures/exp/fig_matrixresgated_sensitivity.pdf}
\caption{Fold-0 MatrixResGated sensitivity scan at \(B=3\). The scan identifies candidate operating regions but is not treated as final hyperparameter evidence.}
\label{fig:sensitivity}
\end{figure}

\subsection{Five-fold check of tuned candidates}

Fold-0 scans are used only for candidate discovery and cannot establish stable performance by themselves. To obtain more reliable evidence, selected MatrixResGated settings were rerun across all five folds. Table~\ref{tab:tuned_candidates} reports these follow-up checks. On DD, the MatrixResGated candidate with hidden dimension 32 reduces parameters from 42,371 to 15,043 while remaining competitive, indicating better parameter efficiency. The lower learning rate is close to the default MatrixResGated benchmark, and the lower gate initialization is weaker. On PROTEINS, the candidate with sparsity strength \(\lambda=0.02\) does not consistently preserve its fold-0 advantage across all folds, so the sparsity gain should not be overstated.

\begin{table}[!htbp]
\centering
\small
\caption{Five-fold checks of selected MatrixResGated candidates.}
\label{tab:tuned_candidates}
\begin{tabular}{l l c c c}
\toprule
Candidate & Dataset & Accuracy & Mean loss & Parameters \\
\midrule
Hidden dimension 32 & DD & \(\mathbf{0.7215\pm0.0330}\) & 0.5717 & 15,043 \\
Gate initialization 0.2 & DD & \(0.7131\pm0.0213\) & 0.5820 & 42,371 \\
Learning rate 0.001 & DD & \(0.7181\pm0.0336\) & 0.5737 & 42,371 \\
Sparsity strength \(\lambda=0.02\) & PROTEINS & \(0.6909\pm0.0447\) & 0.6104 & 38,339 \\
\bottomrule
\end{tabular}
\end{table}

The tuned-candidate check prevents over-claiming from fold-0 sensitivity. It supports a smaller, parameter-efficient MatrixResGated configuration on DD, but it does not support a strong full-fold sparsity claim on PROTEINS.
